\documentclass[10pt,a4paper]{article}
\usepackage{amsmath,amssymb,amsthm,graphicx}
\usepackage{hyperref}
\usepackage{amsmath}
\usepackage{enumerate}
\usepackage{amsfonts}
\usepackage{amsthm}
\newtheorem{theorem}{Theorem}
\newtheorem{proposition}[theorem]{Proposition}
\newtheorem{lemma}[theorem]{Lemma}
\newtheorem{corollary}[theorem]{Corollary}
\newtheorem{rem
ark}[theorem]{Remark}
\newtheorem{definition}[theorem]{Definition}
\newtheorem{example}[theorem]{Example}
\newcommand{\out}{\textup{out}} 
\newcommand{\spa}{\textup{span}} 
\newcommand{\Stab}{\textup{Stab}} 
\newcommand{\peb}{\textup{peb}} 

\newcommand{\AcceptEdge}[1]{\textbf{Accept-Edge}(#1)}
\newcommand{\MovePebble}[1]{\textbf{Move-Pebble}(#1)}
\title{An edge between consecutive vertex classes can be accepted while enough pebbles remain}
\author{Signe Lundqvist, Tovohery Randrianarisoa, Klara Stokes, Joannes Vermant}
\date{arxiv:2306.05050v3; Creative Commons Attribution 4.0 International}
\begin{document}
\maketitle
\noindent\textit{Editorial note.} The counted statement is the author's labelled lemma environment at source0.tex lines 913-917; its proof environment opens at line 918, immediately after \textbackslash{}end\{lemma\}, so the pairing is by adjacency and the retained proof range 918-929 is the whole environment.\par\medskip
\noindent\textit{Editorial note.} Two statements are retained as context: the numbered system of inequalities at source lines 241-245, which the counted statement assumes, and the invariants lemma at source lines 877-885, whose invariant the proof uses at source line 920. These are the only outside labels of the retained material.\par\medskip
\noindent\textit{Editorial note.} The retained proof carries no citation, so no bibliography is needed; the source ships both a bibliography.bib and a prepared retained bibliography for the cards of this source that do cite.\par\medskip
\noindent\textit{Editorial note.} Preamble: the author's own package lines, theorem-environment family and macro block (source lines 4-39); the drawing and algorithm packages of the source are not retained because the retained material does not use them.\par\medskip
\noindent\textit{Editorial note.} Metadata provenance: the licence evidence file records source\_version 2306.05050v3 with status supported and the card uses that value; the author list is transcribed from the authors\_index key of the source record file of the pinned copy, and the title from its title key.\par\medskip
\noindent\textit{Editorial note.} The source carries no label-number dump and no counter values are imposed, so the printed numbering is the wrapper's default; the identity of the counted item is its author label.\par\medskip
\noindent\textit{Editorial scope.} One original result of the pinned author manuscript, counted once: the step of the pebble game accepting an edge between consecutive vertex classes (author label lemma-add-edge(poset)).\par\medskip
\noindent\textit{Allowed standard prerequisites (editorial).}\begin{itemize}
\item The pebble game on a poset. An algorithm that keeps a number of pebbles on each vertex and accepts edges while respecting the sparsity parameters, with a reachability operation used to move pebbles along accepted edges.. It is the algorithm whose behaviour the statement describes.
\item Pebble, span and out counts. For a subset of vertices the number of pebbles, the span and the out-degree are related to the weighted size of the classes the subset meets by the algorithm's invariants.. They are the quantities used to count pebbles inside a reachable set.
\item Reachable set of a vertex. The set of vertices that can be reached from a given vertex by following accepted edges, respecting the direction conventions of the game.. It is the set on which the counting argument is performed.
\end{itemize}\medskip
\begin{align}
    k_i + k_{i+1} &> l_{i, i+1} \label{conditionposet1 -k k l},\\
    k_{i+1} + l_{i-1, i} &> l_{i, i+1}\label{conditionposet2 - k l l},\\
     l_{1,2} \leq  \dots &\leq l_{n-1,n}. \label{conditionposet3-increasing}
\end{align}

\begin{lemma}\label{invariants - posets}
 Let $K\in \mathbb{Z}_{>0}^{n}, L\in \mathbb{Z}_{\geq 0}^{n-1}$ be such that inequalities \eqref{conditionposet1 -k k l}, \eqref{conditionposet2 - k l l} and \eqref{conditionposet3-increasing} are valid. Let $G=(V,E)$ be a graph such that the vertex set has a partition $V=V_1\cup \dots \cup V_n$ such that all edges $e\in E$ have endpoints in $V_{i} \cup V_{i+1}$ for some $1 \leq i < n$. The following are invariants of the algorithm
    \begin{enumerate}
        \item For all $v\in V_i$, one has $\peb(v)+ \out(v) = k_{i}$.
        \item Suppose that $V'\subseteq V$. Then $$\peb(V')+ \spa(V')+ \out(V') = \sum_{i=1}^{n} k_i \vert  V' \cap V_i  \vert.$$
        \item Consider a nonempty subset $V' \subseteq V$. Then $$\peb(V') + \out(V') \geq l(V').$$ 
    \item Consider a subset $V' \subseteq V$. Then $$\spa(V') \leq \sum_{i=1}^{n} k_i \vert  V' \cap V_i \vert - l(V').$$
    \end{enumerate}
\end{lemma}

\begin{lemma}\label{lemma-add-edge(poset)}
 Let $K\in \mathbb{Z}_{>0}^{n}, L\in \mathbb{Z}_{\geq 0}^{n-1}$ be such that inequalities \eqref{conditionposet1 -k k l}, \eqref{conditionposet2 - k l l} and \eqref{conditionposet3-increasing} are valid. Let $G=(V,E)$ be a graph such that the vertex set has a partition $V=V_1\cup \dots \cup V_n$ such that all edges $e\in E$ have endpoints in $V_{i} \cup V_{i+1}$ for some $1 \leq i < n$. 
 
Suppose that after running some steps of the algorithm $D$ is the set of edges we accepted, where all edges in $D$ are between $V_a$ and $V_{a+1}$ for $a\leq i$. Suppose that $e$ is an edge of $G$ which has not been accepted yet such that $e=(uv)$, where $u\in V_i$ and $v\in V_{i+1}$. Then the edge can be accepted (after perhaps moving some pebbles).
\end{lemma}

\begin{proof}
If there are more than $l_{i, i+1}$ pebbles at $u$ and $v$ together, we can add the edge and remove a pebble. So, we may assume that $\peb(u)+\peb(v)<l_{i, i+1}+1$.
Let $V'=\textup{Reach}(u)\cup \textup{Reach}(v)$, and let $E_G(V')$ denote the edges induced by $V'$ in $G$. Note that $V'\cap V_i$ and $V'\cap  V_{i+1}$ are non-empty since $u\in V'\cap V_i$ and $v\in V'\cap V_{i+1}$. Furthermore $V'\cap V_k =\emptyset$, for $k\geq i+2$, since these edges get processed at some later point. Since $G$ is sparse, $$\spa(V')< \vert E_G(V') \vert \leq \sum_{j=1}^{i+1}k_j\vert V' \cap V_j  \vert - l_{i,i+1},$$ where the first inequality is strict since $e\in  E_G(V') \setminus \spa(V')$, and $\max\{l_{k,k+1}\}=l_{i,i+1}$ since the $l_{k,k+1}$ are increasing. By the definition of $\textup{Reach}$, $\out(V')=0$ and by invariant 2 in Lemma~\ref{invariants - posets}, 
\[
\spa(V')+ \peb(V')= \sum_{j=1}^{i+1}k_j\vert V_j \cap V' \vert
\]
and thus
\[
\peb(V') > l_{i, i+1} .
\]
Hence, there is a vertex $w\in V'$ such that $w\notin \{v,u\}$ with at least one pebble left. Since it is reachable by either $v$ or $u$, we can move that pebble (changing the direction of the path). Repeat this process until we get enough pebbles together at $u$ and $v$ so that we can accept the edge $(uv)$ and add it to $D$.
\end{proof}

\end{document}
